\subsection{THE POWER OF R}

We have \(\mathbf{R} = \bigcup_{n\in \mathbf{Z}}[n,n + 1[,\) and all the intervals \([n,n + 1[\) are equipotent to \([0,1[\) . Since \([0,1[\) is an infinite set it follows that \(\mathbf{R}\) is equipotent to the interval \([0,1[\) . By considering the dyadic expansion of the numbers of the interval \([0,1[\) we shall show that this interval is equipotent to the set S of all sequences \((u_n)\) in which each term is equal to 0 or 1.

First, it is equipotent to the subset \(S'\) of \(S\) consisting of sequences \((u_n)\) such that \(u_n = 0\) for an infinity of values of \(n\) . On the other hand, the complement \(S''\) of \(S'\) in \(S\) is equipotent to the set of all improper expansions of rational numbers which are equal to a fraction with denominator \(2^n\) ; these numbers form a subset of \(Q\) , hence the set of them is countable and therefore so is \(S''\) . Since \(S'\) is infinite, it is equipotent to \(S\) , hence the result.

Now S is equipotent to \(\mathfrak{P}(\mathbf{N})\) ; for we can define a bijection of \(\mathfrak{P}(\mathbf{N})\) onto S by mapping each subset A of N to the sequence \((u_n)\) such that \(u_n = 0\) if \(n \in A\) and \(u_n = 1\) if \(n \in \mathbb{C}A\) . We have thus proved:

THEOREM 1 (Cantor). The set of real numbers is equipotent to the set of all subsets of a countably infinite set.

COROLLARY. The set of real numbers has a cardinal strictly greater than the cardinal of a countable set.

A set equipotent to R is said to have the power of the continuum. By Proposition 1 of § 4, no. 1, every interval which contains more than one point has the power of the continuum. Again, the complement of a countable subset of R has the power of the continuum in particular, the set of all irrational numbers has the power of the continuum.

